\mode*
\section{The command model}
\label{sec:shell}

The first educational objective is the shell's command model: the
read--evaluate--print loop, the anatomy of a command line (command,
options, arguments), and the shell as \emph{an} interface to the same
system that the graphical desktop is another interface to.

\subsection{Documented misconceptions and difficulties}

Research on how learners get the command model wrong is scarce
\autocite{Winder2024ShellTutor}; what exists comes largely from
cognitive-psychology studies of Unix users.
Producing correct composite commands is demanding, and help that states
concrete information supports novices where abstract information only helps
experts \autocite{Doane1992UnixPrompts}.
Beginning programming students are commonly introduced through visual
environments rather than the command line \autocite{Dillon2012MentalModels},
so we must assume no prior CLI exposure.
% TODO: The systematic search phase (cf. \cref{sec:method-litreview}) must
% look for newer CS-education sources on CLI misconceptions specifically;
% the Shell Tutor group's logging data may be the closest existing corpus.
% TODO: Add misconceptions observed in our own course's grading and
% tutoring data (datintro), clearly separated from the literature-derived
% ones.

\subsection{Candidate critical aspects}

% XXX Preliminary; to be consolidated once the review phase is complete.
\begin{itemize}
  \item \emph{The command line is a language, not a search box}: the shell
    parses; it does not guess. (Word order, spaces, quoting matter.)
  \item \emph{Textual feedback}: output--and silence--is the system's
    response; error messages are information, not failure.
  \item \emph{The shell acts on the same system as the GUI}: same files,
    same processes---only the interface varies.
\end{itemize}

\subsection{Patterns of variation}

\begin{itemize}
  \item Contrast (interface varies, system invariant): perform one task in
    the GUI, then the same task in the shell; observe the same effect in
    the other interface.
  \item Contrast (one token varies): same command with and without one
    option; with correct and with misspelled argument---observe how the
    response differs.
  \item Generalisation (command invariant, target varies): the same
    \texttt{ls}/\texttt{cp} across different directories and file types.
\end{itemize}
% TODO: Realise these as concrete command sequences (PythonTeX-captured
% terminal transcripts, as the companion papers capture program output).
